\section{Einleitung}
Durch die fortschreitende Digitalisierung und den zunehmenden Ausbau dezentraler 
Energiequellen befinden sich Stromnetze heute in einem großen Wandel. 
Sie entwickeln sich zu Smart Grids. Diese zeichnen sich dadurch aus, 
dass sie physische Energieinfrastruktur mit Informations- und 
Kommunikationstechnologie verbinden. Das soll, trotz schwankender Einspeisung aus 
erneuerbaren Energien und dynamischer Lastprofile, einen stabilen und effizienten 
Netzbetrieb gewährleisten \citep[S.529]{Gungor2013SmartGrid}. Smart Grids unterscheiden 
sich von herkömmlichen Stromnetzen. Letztere waren meist hierarchisch organisiert 
und auf zentrale Energieerzeugung sowie unidirektionale Energieflüsse ausgelegt. 
Im Gegensatz dazu bieten Smart Grids hohe Dezentralität, bidirektionale Energieflüsse,
 Informationsflüsse und adaptive Steuerungsmechanismen \citep[S. 12]{Appelrath2012SmartGrid}.

Diese zunehmende Vernetzung und Abhängigkeiten der Komponenten führt jedoch nicht nur
 zu höherer Effizienz, sondern auch zu einer größeren Verwundbarkeit des Systems. 
 Störungen durch zufällige Ausfälle, wie Überlastungen oder auch gezielte Angriffe 
 auf bestimmte Knotenpunkte, können sich aufgrund der wechselseitigen 
 Interdependenzen schnell zu großflächigen Ausfällen ausbreiten und zu 
 Versorgungsstörungen im ganzen Netz entwickeln \citep[S. 1025]{Buldyrev2010Interdependent}. 
 Gerade die enge Zusammenarbeit von physischen und digitalen Netzen macht Smart Grids 
 anfällig für kaskadierende Effekte, die mit Analysen von konventionellen Stromnetzen 
 nicht ausreichend erfasst werden können.
Vor diesem Hintergrund gewinnt der Begriff der Resilienz ganz besonders an Bedeutung,
 denn er beschreibt die Fähigkeit eines Systems Störungen zu absorbieren und seine grundlegende Funktion aufrechtzuerhalten oder sich nach einem Ausfall wieder zu erholen \citep[S. 17]{Holling1973Resilience}. Zur systematischen Analyse dieser Eigenschaften bietet die Theorie komplexer Netzwerke einen geeigneten methodischen Rahmen.
Für die Fehlertoleranz und Robustheit des Systems spielt dabei die Struktur des Netzwerks eine wichtige Rolle. Denn unterschiedliche Topologien, wie zufällige Netzwerke, Smallworld- oder Skalenfreie Netzwerke, weisen unterschiedliche charakteristische Muster auf und gehen unterschiedlich mit zufälligen Ausfällen und gezielten Angriffen um. Insbesondere unterscheiden sich die Rollen hochvernetzter Knoten und die Fragmentierung der Netzwerke \citep[S. 105ff]{Albert2000ErrorAttack}.
Ausgehend davon verfolgt diese Arbeit das Ziel, die folgende Leitfrage zu untersuchen: „Wie beeinflusst die Netzwerkstruktur eines Smart Grids die Fehlertoleranz beziehungsweise die Robustheit gegenüber zufälligen Ausfällen und gezielten Angriffen?" Mithilfe netzwerkbasierter Methoden werden dazu unterschiedliche Störszenarien simuliert und geeignete Resilienzmetriken berechnet, um strukturelle Einflussfaktoren zu identifizieren und Ansatzpunkte zur gezielten Verbesserung der Netzstabilität aufzuzeigen.





